\section{Consequences}\label{sec:consequences}

\subsection{The elementary expansion}

\begin{proof}[Proof of Corollary~\ref{cor:circularExpansion}]
By Lemma~\ref{lem:stateExpansion}, and since $\asc$ is constant on flip classes,
\[
 \LLT_\Gamma(\xvec;q+1)=\sum_{\theta}q^{\asc(\theta)}\Kp_{\theta}(\xvec)
 =\sum_Cq^{\asc(C)}\Kv(C),
\]
the second sum ranging over the flip classes.
By Theorem~\ref{thm:main} this equals
\[
 \sum_Cq^{\asc(C)}\Ev(C)=\sum_\theta q^{\asc(\theta)}\elementaryE_{\lambda(\theta)}(\xvec).
\qedhere
\]
\end{proof}

Thus the coefficient of $q^k\elementaryE_\lambda$ in $\LLT_\Gamma(\xvec;q+1)$ is the
number of admissible orientations $\theta$ with $\asc(\theta)=k$ and
$\lambda(\theta)=\lambda$, and every flip class contributes a sum of such terms.
Substituting $q-1$ for $q$, we also have
$\LLT_\Gamma(\xvec;q)=\sum_\theta(q-1)^{\asc(\theta)}\elementaryE_{\lambda(\theta)}(\xvec)$,
which is the form in which the expansion was stated
in~\cite{AlexanderssonSulzgruber2022}.
For ordinary graphs, Corollary~\ref{cor:circularExpansion} is
Theorem~\ref{thm:AS}, and the proof above does not use the recursion
of~\cite{AlexanderssonSulzgruber2022}.
For circular graphs it answers the question
of~\cite[Sec.~6.8]{AlexanderssonSulzgruber2022}.

\subsection{LLT cumulants}

We now refine the elementary expansion of LLT cumulants of M.~Dołęga and
M.~Kowalski~\cite{DolegaKowalski2021x} to flip classes, for circular graphs.
Let $\Gamma$ be a circular decorated unit-interval graph and $U$ a nonempty set of
residues.
List the lifts of the elements of $U$ increasingly as $(u_j)_{j\in\setZ}$, so that
$u_{j+m}=u_j+n$ with $m=|U|$, and put $b'(j)=\min\{k:u_k\geq b(u_j)\}$.
Then $b'$ is weakly increasing, $j-m<b'(j)\leq j$, and $b'(j+m)=b'(j)+m$, so it
defines a circular unit-interval graph on $U$ whose lifted natural edges are the
pairs $u_ku_j$ with $b(u_j)\leq u_k<u_j$.
A strict edge $u_ku_j$ of $\Gamma$ remains a corner, since $b(u_j)=u_k<b(u_j+1)$ gives
$b'(j)=k<b'(j+1)$.
We call the resulting graph the \defin{induced} graph $\Gamma[U]$.
Its natural edges are exactly the natural edges of $\Gamma$ with both endpoints in
$U$, and its strict edges are the strict edges of $\Gamma$ inside $U$.

Let $r\geq1$, and let $f:\setZ/n\setZ\to[r]$ be surjective and satisfy
$f(u)=f(v)$ for every strict edge $uv$.
For nonempty $B\subseteq[r]$ write $\Gamma[B]$ for the graph induced on
$f^{-1}(B)$.
Let $\Pi_r$ be the lattice of set partitions of $[r]$, with maximum $\hat1=\{[r]\}$.
Following~\cite{DolegaKowalski2021x}, the \defin{LLT cumulant} is
\[
 \defin{\kappa_{\Gamma,f}(\xvec;q)}\coloneqq(q-1)^{1-r}\sum_{\pi\in\Pi_r}
 (-1)^{|\pi|-1}(|\pi|-1)!\prod_{B\in\pi}\LLT_{\Gamma[B]}(\xvec;q).
\]
For an admissible orientation $\theta$, let $\sigma_f(\theta)$ be the set partition of
$[r]$ into the connected components of the undirected multigraph on $[r]$ with an
edge $f(u)f(v)$ for every edge $uv$ of $\theta\cup S$.
We call $\theta$ \defin{$f$-connected} if $\sigma_f(\theta)=\hat1$.

\begin{proposition}[{cf.~\cite[Thm.~3.8]{DolegaKowalski2021x}}]\label{prop:cumulant}
We have
\[
 \kappa_{\Gamma,f}(\xvec;q+1)=\sum_{\theta}q^{\asc(\theta)+1-r}\,\Kp_{\theta}(\xvec),
\]
the sum ranging over the $f$-connected admissible orientations of $\Gamma$, and all
exponents are nonnegative.
\end{proposition}

\begin{proof}
For $\pi\in\Pi_r$, the orientations of $\Gamma$ without ascending edges between
the sets $f^{-1}(B)$, $B\in\pi$, are the tuples of orientations of the graphs
$\Gamma[B]$, since no strict edge joins two such sets.
By Lemma~\ref{lem:stateExpansion} for each $\Gamma[B]$,
\[
 \prod_{B\in\pi}\LLT_{\Gamma[B]}(\xvec;q+1)
 =\sum_{\sigma_f(\theta)\leq\pi}q^{\asc(\theta)}\Kp_{\theta}(\xvec).
\]
Inserting this in the definition, the coefficient of a fixed $\theta$ is
$q^{\asc(\theta)+1-r}$ times
$\sum_{\sigma_f(\theta)\leq\pi}(-1)^{|\pi|-1}(|\pi|-1)!$.
Since $(-1)^{|\pi|-1}(|\pi|-1)!$ is the Möbius function of $\Pi_r$ from $\pi$ to
$\hat1$, this sum is $[\sigma_f(\theta)=\hat1]$.
Since $f(u)=f(v)$ for every strict edge $uv$, strict edges become loops in the
multigraph on $[r]$.
They do not connect different vertices, so an $f$-connected orientation has at
least $r-1$ ascending edges.
\end{proof}

By Lemma~\ref{lem:flipInvariants}, being $f$-connected is constant on flip classes.
Grouping Proposition~\ref{prop:cumulant} by flip classes and applying
Theorem~\ref{thm:main}, we obtain the following.

\begin{corollary}\label{cor:cumulant}
With the notation above,
\[
 \kappa_{\Gamma,f}(\xvec;q+1)=\sum_Cq^{\asc(C)+1-r}\sum_{\theta\in C}
 \elementaryE_{\lambda(\theta)}(\xvec),
\]
where $C$ ranges over the flip classes of $f$-connected admissible orientations.
In particular, every flip class contributes an elementary-positive summand.
\end{corollary}

For ordinary graphs, summing over the classes recovers the elementary expansion
of~\cite[Thm.~3.17]{DolegaKowalski2021x}.

\subsection{Questions}

For a unit-interval graph without strict edges, the chromatic quasisymmetric
function and the unicellular LLT polynomial are related by a
plethysm, see~\cite[Lem.~6.17]{AlexanderssonSulzgruber2022}.
OpenAI proved that the elementary coefficients of the full chromatic
quasisymmetric function belong to
$\mathbb{N}[q]$~\cite[Thm.~1.1]{OpenAI2026}.
B.-H.~Hwang's stronger classwise conjecture remains open: each flip class of
acyclic orientations should contribute an elementary-positive
summand~\cite[Conj.~3.10]{Hwang2022x}.
Theorem~\ref{thm:main} therefore establishes the LLT classwise analogue before
its chromatic counterpart.

The two kinds of flip give different components.
Hwang reverses a source or sink fork whose outer vertices are nonadjacent,
whereas Definition~\ref{def:flip} reverses the two short edges of a triangle and
retains its long edge.
Moreover, Hwang's classes consist of acyclic orientations, while our classes
range over all orientations of the normal edges.
The plethystic relation does not preserve elementary positivity term by term:
already $(q-1)^{-1}\elementaryE_2[(q-1)\xvec]=(q+1)\elementaryE_2-\elementaryE_{11}$.

\begin{question}
Is there a direct relation between Hwang's flip classes and the flip classes of
Definition~\ref{def:flip} which transfers classwise elementary positivity
between the chromatic and LLT settings?
\end{question}

\begin{question}
Give a bijective proof of Theorem~\ref{thm:main}, that is, a bijection between the
strict $\Pplus{\theta}$-partitions for $\theta\in C$ and the strict partitions of the
chains of the highest reachable fibres for $\theta\in C$, preserving the multiset of
colors.
\end{question}
